% EJEMPLO OPCIONAL. Insertar desde main.tex:
% % EJEMPLO OPCIONAL. Insertar desde main.tex:
% % EJEMPLO OPCIONAL. Insertar desde main.tex:
% % EJEMPLO OPCIONAL. Insertar desde main.tex:
% \input{ejemplos/matriz_operacionalizacion.tex}
%
% Los datos son ficticios y deben reemplazarse. No colocar este ejemplo
% en un informe final sin verificar variables, dimensiones e indicadores.
\begin{matrizOperacionalizacion}
\noindent\textbf{Matriz de operacionalización de variables}\par\vspace{6pt}
\tituloMatrizOperacionalizacion[tab:operacionalizacion]{Operacionalización de las variables de estudio}
\iniciotablaHorizontal
\noindent\begin{tabularx}{\linewidth}{|
>{\raggedright\arraybackslash}p{2.15cm}|
>{\raggedright\arraybackslash}p{3.1cm}|
>{\raggedright\arraybackslash}p{3cm}|
>{\raggedright\arraybackslash}X|
>{\raggedright\arraybackslash}p{3.9cm}|
>{\raggedright\arraybackslash}p{1.7cm}|}\hline
\textbf{Variable}&
\textbf{Definición operacional}&
\textbf{Dimensión}&
\textbf{Indicador}&
\textbf{Fuentes / instrumentos}&
\textbf{Escala}\\\hline
\multirow{2}{2.15cm}{\raggedright Variable X: [nombre]}&
\multirow{2}{3.1cm}{\raggedright [Definición medible de la variable X]}&
X1. [Primera dimensión]&
\celdaIndicadores{\item \relax [Primer indicador de X1].\item \relax [Segundo indicador de X1].}&
[Fuentes o instrumentos]&
[Razón, ordinal, etc.]\\\cline{3-6}
&&X2. [Segunda dimensión]&
\celdaIndicadores{\item \relax [Primer indicador de X2].\item \relax [Segundo indicador de X2].}&
[Fuentes o instrumentos]&
[Escala]\\\hline
\multirow{2}{2.15cm}{\raggedright Variable Y: [nombre]}&
\multirow{2}{3.1cm}{\raggedright [Definición medible de la variable Y]}&
Y1. [Primera dimensión]&
\celdaIndicadores{\item \relax [Primer indicador de Y1].\item \relax [Segundo indicador de Y1].}&
[Fuentes o instrumentos]&
[Escala]\\\cline{3-6}
&&Y2. [Segunda dimensión]&
\celdaIndicadores{\item \relax [Primer indicador de Y2].\item \relax [Segundo indicador de Y2].}&
[Fuentes o instrumentos]&
[Escala]\\\hline
\end{tabularx}
\fuenteTabla{Elaboración propia. Cada indicador se presenta individualmente dentro de su dimensión.}
\end{matrizOperacionalizacion}

%
% Los datos son ficticios y deben reemplazarse. No colocar este ejemplo
% en un informe final sin verificar variables, dimensiones e indicadores.
\begin{matrizOperacionalizacion}
\noindent\textbf{Matriz de operacionalización de variables}\par\vspace{6pt}
\tituloMatrizOperacionalizacion[tab:operacionalizacion]{Operacionalización de las variables de estudio}
\iniciotablaHorizontal
\noindent\begin{tabularx}{\linewidth}{|
>{\raggedright\arraybackslash}p{2.15cm}|
>{\raggedright\arraybackslash}p{3.1cm}|
>{\raggedright\arraybackslash}p{3cm}|
>{\raggedright\arraybackslash}X|
>{\raggedright\arraybackslash}p{3.9cm}|
>{\raggedright\arraybackslash}p{1.7cm}|}\hline
\textbf{Variable}&
\textbf{Definición operacional}&
\textbf{Dimensión}&
\textbf{Indicador}&
\textbf{Fuentes / instrumentos}&
\textbf{Escala}\\\hline
\multirow{2}{2.15cm}{\raggedright Variable X: [nombre]}&
\multirow{2}{3.1cm}{\raggedright [Definición medible de la variable X]}&
X1. [Primera dimensión]&
\celdaIndicadores{\item \relax [Primer indicador de X1].\item \relax [Segundo indicador de X1].}&
[Fuentes o instrumentos]&
[Razón, ordinal, etc.]\\\cline{3-6}
&&X2. [Segunda dimensión]&
\celdaIndicadores{\item \relax [Primer indicador de X2].\item \relax [Segundo indicador de X2].}&
[Fuentes o instrumentos]&
[Escala]\\\hline
\multirow{2}{2.15cm}{\raggedright Variable Y: [nombre]}&
\multirow{2}{3.1cm}{\raggedright [Definición medible de la variable Y]}&
Y1. [Primera dimensión]&
\celdaIndicadores{\item \relax [Primer indicador de Y1].\item \relax [Segundo indicador de Y1].}&
[Fuentes o instrumentos]&
[Escala]\\\cline{3-6}
&&Y2. [Segunda dimensión]&
\celdaIndicadores{\item \relax [Primer indicador de Y2].\item \relax [Segundo indicador de Y2].}&
[Fuentes o instrumentos]&
[Escala]\\\hline
\end{tabularx}
\fuenteTabla{Elaboración propia. Cada indicador se presenta individualmente dentro de su dimensión.}
\end{matrizOperacionalizacion}

%
% Los datos son ficticios y deben reemplazarse. No colocar este ejemplo
% en un informe final sin verificar variables, dimensiones e indicadores.
\begin{matrizOperacionalizacion}
\noindent\textbf{Matriz de operacionalización de variables}\par\vspace{6pt}
\tituloMatrizOperacionalizacion[tab:operacionalizacion]{Operacionalización de las variables de estudio}
\iniciotablaHorizontal
\noindent\begin{tabularx}{\linewidth}{|
>{\raggedright\arraybackslash}p{2.15cm}|
>{\raggedright\arraybackslash}p{3.1cm}|
>{\raggedright\arraybackslash}p{3cm}|
>{\raggedright\arraybackslash}X|
>{\raggedright\arraybackslash}p{3.9cm}|
>{\raggedright\arraybackslash}p{1.7cm}|}\hline
\textbf{Variable}&
\textbf{Definición operacional}&
\textbf{Dimensión}&
\textbf{Indicador}&
\textbf{Fuentes / instrumentos}&
\textbf{Escala}\\\hline
\multirow{2}{2.15cm}{\raggedright Variable X: [nombre]}&
\multirow{2}{3.1cm}{\raggedright [Definición medible de la variable X]}&
X1. [Primera dimensión]&
\celdaIndicadores{\item \relax [Primer indicador de X1].\item \relax [Segundo indicador de X1].}&
[Fuentes o instrumentos]&
[Razón, ordinal, etc.]\\\cline{3-6}
&&X2. [Segunda dimensión]&
\celdaIndicadores{\item \relax [Primer indicador de X2].\item \relax [Segundo indicador de X2].}&
[Fuentes o instrumentos]&
[Escala]\\\hline
\multirow{2}{2.15cm}{\raggedright Variable Y: [nombre]}&
\multirow{2}{3.1cm}{\raggedright [Definición medible de la variable Y]}&
Y1. [Primera dimensión]&
\celdaIndicadores{\item \relax [Primer indicador de Y1].\item \relax [Segundo indicador de Y1].}&
[Fuentes o instrumentos]&
[Escala]\\\cline{3-6}
&&Y2. [Segunda dimensión]&
\celdaIndicadores{\item \relax [Primer indicador de Y2].\item \relax [Segundo indicador de Y2].}&
[Fuentes o instrumentos]&
[Escala]\\\hline
\end{tabularx}
\fuenteTabla{Elaboración propia. Cada indicador se presenta individualmente dentro de su dimensión.}
\end{matrizOperacionalizacion}

%
% Los datos son ficticios y deben reemplazarse. No colocar este ejemplo
% en un informe final sin verificar variables, dimensiones e indicadores.
\begin{matrizOperacionalizacion}
\noindent\textbf{Matriz de operacionalización de variables}\par\vspace{6pt}
\tituloMatrizOperacionalizacion[tab:operacionalizacion]{Operacionalización de las variables de estudio}
\iniciotablaHorizontal
\noindent\begin{tabularx}{\linewidth}{|
>{\raggedright\arraybackslash}p{2.15cm}|
>{\raggedright\arraybackslash}p{3.1cm}|
>{\raggedright\arraybackslash}p{3cm}|
>{\raggedright\arraybackslash}X|
>{\raggedright\arraybackslash}p{3.9cm}|
>{\raggedright\arraybackslash}p{1.7cm}|}\hline
\textbf{Variable}&
\textbf{Definición operacional}&
\textbf{Dimensión}&
\textbf{Indicador}&
\textbf{Fuentes / instrumentos}&
\textbf{Escala}\\\hline
\multirow{2}{2.15cm}{\raggedright Variable X: [nombre]}&
\multirow{2}{3.1cm}{\raggedright [Definición medible de la variable X]}&
X1. [Primera dimensión]&
\celdaIndicadores{\item \relax [Primer indicador de X1].\item \relax [Segundo indicador de X1].}&
[Fuentes o instrumentos]&
[Razón, ordinal, etc.]\\\cline{3-6}
&&X2. [Segunda dimensión]&
\celdaIndicadores{\item \relax [Primer indicador de X2].\item \relax [Segundo indicador de X2].}&
[Fuentes o instrumentos]&
[Escala]\\\hline
\multirow{2}{2.15cm}{\raggedright Variable Y: [nombre]}&
\multirow{2}{3.1cm}{\raggedright [Definición medible de la variable Y]}&
Y1. [Primera dimensión]&
\celdaIndicadores{\item \relax [Primer indicador de Y1].\item \relax [Segundo indicador de Y1].}&
[Fuentes o instrumentos]&
[Escala]\\\cline{3-6}
&&Y2. [Segunda dimensión]&
\celdaIndicadores{\item \relax [Primer indicador de Y2].\item \relax [Segundo indicador de Y2].}&
[Fuentes o instrumentos]&
[Escala]\\\hline
\end{tabularx}
\fuenteTabla{Elaboración propia. Cada indicador se presenta individualmente dentro de su dimensión.}
\end{matrizOperacionalizacion}
